\section{Introduction}\label{sec:intro}

Multi-object tracking (MOT) finds every object of interest in a video and gives each one an
identity number (ID) that stays the same over time. It is used in surveillance, driving, and drone
systems \cite{zhu2022visdrone}, where results are needed with little delay. Most
MOT systems follow the tracking-by-detection design \cite{adzemovic2026survey}: in each frame, a
detector returns candidate boxes with confidence scores, and a tracker links them into tracks, one
per object \cite{bewley2016sort,zhang2022bytetrack,cao2023ocsort}. The tracker uses a box only if
its score passes a fixed threshold. A two-stage tracker, such as ByteTrack, also uses lower-score
boxes in a second step to continue unmatched tracks \cite{zhang2022bytetrack}; a single-stage
tracker, such as OC-SORT, ignores every box below one threshold \cite{cao2023ocsort}.

In real systems the detector is often replaced, retrained, or compressed, while the tracker keeps
the thresholds that its authors set for another detector. Detector scores are not calibrated across
models and conditions \cite{kuppers2020detcalib}, so a new detector usually gives scores on a
different scale, and nothing in the system notices. The effect can be large: halving the scores of
a published detector reduces two well-known trackers to 0 higher order tracking accuracy (HOTA)
\cite{luiten2021hota}, because no detection passes their thresholds (Sect.~\ref{sec:shift}). For example, ByteTrack in
its official setting starts a track only from a box with a score of at least 0.7; a car that the
original detector scored 0.9 now gets 0.45, so ByteTrack never starts a track, and its HOTA falls
from 67.70 to 0. \scl{} finds from the scores themselves that the confident group now begins
lower and lets it reach the tracker again, and ByteTrack reaches 66.17 HOTA again, with no labels and no change to its code.
Retuning every tracker for every detector needs labels, which a working system usually does not
have. Adaptive-threshold methods are each built into one tracker
\cite{ma2023adaptivebytetrack,shim2024adaptrack,act2026adaptive}, score calibration needs
labelled data \cite{kuppers2020detcalib}, and adaptive inference methods learn their settings for
one model \cite{han2022dynamic}.

This paper presents a self-calibrating control layer (\scl{}) that sits between a detector and a
tracker that are both kept fixed (frozen): their code, weights, and stored settings are not
changed. We call the tracker that \scl{} serves the \emph{host}. The novelty of \scl{} is
\emph{selective, host-aware adaptation}. It is selective because it first decides, from the
detector scores alone and without labels, whether the scores of the stream fit the host at all:
when most higher scores of recent frames form a clearly confident group, the stream is clean, and
the host keeps its scores and, unless they would reject the confident group, its thresholds. Only
when many scores are uncertain does \scl{} take control: it chooses which candidates reach the
host, gives them new scores, and sets the host's thresholds for the frame. It is host-aware because
every action is expressed in the host's own terms: four published settings of the host, the host
contract (Sect.~\ref{sec:contract}), tell the layer where the host's thresholds lie and whether it
has a low-score step, so the same frozen rules serve single-stage and two-stage trackers.

With halved detector scores, \scl{} restored two trackers from 0 to about 66 HOTA; with other
detectors, it raised the multiple object tracking accuracy (MOTA) of two-stage trackers by up to
24.7 points and cut their identity switches by 43\% to 56\%, while well-matched strong trackers
stayed within 0.06 HOTA points of their own accuracy.

The main contributions are:
\begin{enumerate}
\item \emph{Selective adaptation.} A label-free, causal test of detector--tracker compatibility:
      nested Otsu splits of recent logits and a whole-stream majority rule decide whether to act,
      so a well-matched system is left as it is and a mismatched one is corrected
      (Sects.~\ref{sec:calib}--\ref{sec:choice}).
\item \emph{Host-aware adaptation.} A host contract of four published values through which one
      fixed configuration sets the thresholds, scores, and matching tolerance of unchanged
      single-stage and two-stage trackers, with a track-guided recovery of weak detections for
      trackers that lack a low-score step (Sects.~\ref{sec:rescue}--\ref{sec:contract}).
\item \emph{Evidence for compatibility.} An evaluation of one frozen configuration on seven
      trackers, four detection sources, and three datasets, with paired bootstrap intervals,
      comparisons with adaptive thresholds and the published ACT method, a component ablation and a
      parameter sensitivity analysis (Sect.~\ref{sec:results}), the recorded development path
      (Sect.~\ref{sec:development}), and a cost of about 3 ms per frame on an ordinary processor
      (central processing unit, CPU) (Sect.~\ref{sec:discussion}).
\end{enumerate}

Section~\ref{sec:related} covers the data, metrics, and related work;
Sections~\ref{sec:method} and~\ref{sec:development} present \scl{} and its development;
Section~\ref{sec:results} gives the results; and Sections~\ref{sec:discussion}
and~\ref{sec:conclusion} discuss and conclude.
